\section{Survey of UE for molecular force fields}\label{sec:survey}
\begin{table*}
    \centering
    \resizebox{\linewidth}{!}{
    \definecolor{mygreen}{HTML}{3C8031}
\definecolor{myred}{HTML}{A9341F}

\newcommand{\tick}{{\color{mygreen}\cmark}}
\newcommand{\cross}{{\color{red}\xmark}}
\newcommand{\neutral}{{\color{orange}$\bm{\sim}$}}

\begin{tabular}{lcccccccc}
\toprule
{} & Ensemble & MC-Dropout & Evidential & GP & SVGP & SVGP-DKL & \oursacro{}\\
\midrule
\desone{} & \tick & \tick & \tick & \tick & \cross & \tick & \tick \\
\destwo{} & \tick & \tick & \tick & \tick & \cross & \tick & \tick \\
\desfour{} &  \tick & \neutral & \cross & \cross & \cross & \cross & \tick  \\
\desfive{} & \tick & \neutral & \tick & \cross & \cross & \neutral & \tick \\
\dessix{} & \tick & \tick & \neutral & \tick & \tick & \tick & \tick \\
\desseven{} & \cross & \cross & \tick & \cross & \tick & \tick & \tick \\
\bottomrule
\end{tabular}

    }
    \caption{Desiderata fulfillment of different uncertainty estimation methods.}
    \label{tab:comp_methods}
\end{table*}
While there is a variety of UE methods for deep learning\footnote{we refer the interested reader to \citet{gawlikowski2021survey}}, here we only discuss (twice) differentiable solutions that fit the application of predicting molecular energies and forces.
UE methods can be broadly categorized into two families of methods: \emph{sampling-based} and \emph{sampling-free} methods.
For each method, we outline its advantages and disadvantages and analyze it based on our desiderata set in \autoref{sec:desiderata}.
An overview of the methods can be found in \autoref{tab:comp_methods}.


\textbf{Ensembles} are touted as the `gold-standard' solution when it comes to UE.
An ensemble is a collection of different methods trained on the same dataset \cite{dataset-shift}.
By varying between architecture or initialization between the individual models, one obtains several estimates about the desired property.
By computing statistical metrics of these samples, e.g., their standard deviation, one obtains an uncertainty quantification \cite{ensembles}. In mathematical notation this can be expressed as: 
\begin{equation}\label{eq:ensemble_uncertainty_estimate}\begin{aligned}
    &\uncertainty_{\energy} = \sqrt{\variance \left(f_1(\molecule), \dots, f_N(\molecule)\right)}\\
    &\uncertainty_{\forces} = \trace \left[\covariance\left(\frac{\partial f_1(\molecule)}{\partial \pos}, \dots \frac{\partial f_N(\molecule)}{\partial \pos} \right) \right].
\end{aligned}\end{equation}
Here $u_{\energy}, u_{\forces}$ are the measures of the predictive uncertainty for the energy and forces. $f_i$ denotes the i-th component of the ensemble, i.e., the i-th trained network. We use the trace as a metric to capture the uncertainty of the forces covariance matrix in a single number.
In terms of our desiderata, ensembles generally preserve the properties of the underlying methods.
Thus, they satisfy the same desiderata as their components with \desseven{} being a notable exception.
As one must train several models \emph{and} perform inference multiple times, an ensemble's runtime scales linearly with its size. 
In case one uses ensembles of GNNs~\citep{gasteigerDirectionalMessagePassing2019} one directly satisfies the physics-informed. 
As we show in \autoref{sec:experiments}, ensembles also empirically perform well on \desfive{} and \dessix{} but crucially fail at \desseven{}.

\textbf{Monte Carlo Dropout} (MCD) is proposed as a comparatively efficient alternative to ensembles with the idea of merging the ensemble into a single model with probabilistic output. Instead of having different models as in ensembles, we leverage the dropout during inference and get multiple predictions for the same molecule. Hence, the formulation for $u_{\energy}, u_{\forces}$ is the same as in \autoref{eq:ensemble_uncertainty_estimate} while the only difference is that $f_i$ denotes a single evaluation of the same model but with different dropped connections \cite{gal2016dropout}. As one simply drops neurons within one model, only one model has to be trained.
Like ensembles, MCD generally inherits the desiderata from its base model but only in expectation as it randomly drops connections.
As we show in \autoref{sec:experiments}, MCD deteriorates the predictive performance.
Further, we observe that while in expectation locality preserving, MCD's uncertainty is not behaving size-consistent for practical sample sizes and is thus not fulfilling \desfour{}.
Lastly, while one only has to train a single model compared to ensembles, one still must perform multiple passes during inference violating the \desseven{} desiderata.


\textbf{Evidential Regression}:
In evidential regression, one directly parametrizes an evidential distribution over the target rather than the target~\citep{soleimany2021evidential}.
In practice, this means that instead of directly estimating a normal distribution on the energy, the model estimates the parameters of an evidential Normal-Inverse-Gamma distribution \cite{evidential-regression,charpentier_natural_2022}. 
Given this distribution, one can perform statistical tests to quantify uncertainties.
While only requiring a single forward pass, this approach does not necessarily yield trustworthy uncertainty estimates, as we will see in \autoref{sec:experiments}.
Similar to the unreliability of the energy prediction itself, the outputted distribution parameters are similarly subject to unpredictable behavior outside the training regime \cite{charpentierPosteriorNetworkUncertainty2020, dbu-robustness}.
Further, by predicting a single distribution for the entire molecular graph, evidential regression requires global embeddings violating \desfour{}.

\begin{figure*}[t!]
    \centering
    \includegraphics[width=\textwidth]{figures/model_figure_v4.pdf}
    \vspace{-0.15in}
    \caption{Overview of \ours{}}
    \label{fig:model_figure}
    \vspace{-0.15in}
\end{figure*}

\textbf{Gaussian Processes} (GPs) are a type of non-parametric models for distributions over functions that are known for their flexibility and expressiveness \cite{bishop, Jakkala2021DeepGP}. 
They are particularly frequently used in UE.
One of the key advantages of using GPs is that they allow for the modeling of complex relationships between variables, while also providing a way to quantify the uncertainty associated with the predictions made by the model.
In the following, we adapt the notation from \citet{Rasmussen2009GaussianPF} to our setting.
Given the data $\dataset$, we assume a GP prior: $f(\dataset) \sim \gp(\mu(\dataset), k_\gamma(\dataset, \dataset^{\transpose}))$, where $\mu(\cdot)$ is the mean function and $k_\gamma(\cdot, \cdot)$ is the covariance function with parameters $\gamma$. Any collection of points $\pfun$ lying on $f(\dataset)$ follows a normal distribution as:
\begin{equation}
\begin{aligned}
	p(\pfun) &= p(f(\dataset)) \\&= p(\left[f(\dataset_1), \dots, f(\dataset_n) \right]^\top)  \sim \normal (\pfun \vert\, \svect \mu, \matr K_{\dataset \dataset}),
 \end{aligned}
\end{equation}
where $\svect \mu_i = \mu(\dataset_i)$ is the mean vector and $\matr K^{\dataset, \dataset}_{ij} = k_\gamma(\dataset_i, \dataset_j)$ the Gram matrix. Assuming that our labels are given as noisy observations of the latent variable, we model the prediction as: $\vect y_i = f(\dataset_i) + \epsilon$, where $\epsilon \sim \normal(0, \sigma)$. We can then predict the values for a new collection of inputs, $\dataset_\star$, via the predictive distribution:
\begin{equation}
	p(\pfun_\star \,\vert \, \dataset_\star, \dataset, \vect y, \sigma, \gamma) \sim \normal (\pfun_\star \vert\, \EX [\pfun_\star], \Cov(\pfun_\star))
\end{equation}
\begin{equation}
	\EX [\vect {f}_\star] = \mu_{\dataset_\star} + \matr K_{\dataset_\star \dataset} \left[K_{\dataset \dataset} + \sigma ^2 \matr I \right]^{-1} \vect y, \quad 
\end{equation}
\begin{equation}
	\Cov(\pfun_\star) =  \matr K_{\dataset_\star \dataset_\star} - \matr K_{\dataset_\star \dataset}\left[\matr K_{\dataset \dataset} + \sigma ^2 \matr I\right] ^{-1}\matr K_{\dataset \dataset_\star}.
\end{equation}
As shown by \citet{vandermause2020On}, we can define a GP acting on the input and hand-crafted features to predict the forces of a molecule with an energy-conserving model while being rotationally covariant. However, due to the inversion of the Gram matrix, the exact GP inference scales with $\mathcal{O}(\ndatapoints^3)$ where $\ndatapoints$ is the number of columns or rows of $\dataset$. Furthermore, compared to recently proposed GNNs the GPs only achieve low accuracies~\citep{gasteiger2021gemnet,batznerEquivariantGraphNeural2022}. 

The desiderata \desfive{} and \desseven{} are thus not fulfilled. 
Further, because a GP compares the graphs as a whole, it does not scale additively with system size, thereby violating \desfour{}.

\textbf{Stochastic Variational Gaussian Process} (SVGP): As the exact GP requires the inversion of the covariance matrix, it comes with a significant computational burden. To mitigate this issue, sparse approximations and variational inference have been proposed, known as Sparse and Variational GP (SVGP) \citep{titsias2009Variational, hensman2013Gaussian, hensman2014Scalable}. 

The SVGP approach entails the selection of a subset of data points, referred to as inducing points $\inducingpos$ and variables $\inducingvars$ to approximate the underlying distribution through the use of variational inference. A prior can be imposed on the joint distribution:\footnote{We adapt \citet{hensman2014Scalable}'s notation for consistency.}
\begin{equation}
	p(\pfun, \inducingvars) \sim \normal \left(
	\begin{aligned}
		\left[
			{\begin{array}{c}
				\pfun\\
				\inducingvars\\
			 \end{array}}
		\right] 
	\end{aligned}\vert\; \vect 0, 
	\begin{aligned}
		\left[
			{\begin{array}{cc}
				\matr K_{\dataset \dataset} & \matr K_{\inducingpos\dataset}\\
				\matr K_{\inducingpos \dataset} & \matr K_{\inducingpos \inducingpos}\\
			\end{array}}
		\right]	
	\end{aligned}
	\right).
\end{equation} 
Note, that this equation describing the relationship between latent variables $\pfun$ and inducing variables $\inducingvars$ is equivalent to the relationship between the training data and the new test points for an exact GP. 

To learn the inducing points and variables, we need to include these into the optimization formulation. \citet{titsias2009Variational} introduced variational inference to approximate the posterior $p(f, u \vert \vect y)$ by $q(\pfun, \inducingvars)$. With the introduction of additional variational parameters, \citet{hensman2013Gaussian} extended this approach and enabled stochastic gradient optimization. Assuming $q(\pfun, \inducingvars) = p(\pfun \vert \inducingvars) q(\inducingvars)$ and $q(\inducingvars)$ being gaussian, we can lower bound the marginal likelihood by the Evidence Lower Bound (ELBO):
\begin{equation}\begin{aligned}\label{eq:elbo}
	\log &p(\pfun_\star) \\ &\geq \EX_{q(\inducingvars)p(\pfun\vert \inducingvars)} \left[\log p(\pfun_\star \vert \pfun)\right] - \KL\left(q(\inducingvars) \vert\vert p(\inducingvars)\right).
\end{aligned}\end{equation}
This ELBO can be used as loss function for stochastic gradient optimization. The computational complexity of an SVGP is in $\mathcal{O}(\ninducingpoints^2\ndatapoints)$ where $m$ is the number of inducing points. Important to note, is that the inducing points and variables are optimized to best fit the distribution and are not instances of the training data. 

Compared to the exact GP, the SVGP fulfills the \desseven{} desideratum. However, as the inducing points are free variables to be optimized, their location information has no connection to the atom positions breaking rotational invariance~\citep{vandermause2020On}).
Thus, compared to GPs, \desone{} is not fulfilled anymore. 
